%=======================================================================
%  Gabriel Adrian Manas — Data Analyst / Data Scientist resume
%  Tailored from the Medical Physics CV toward Data Science / Analytics.
%  Compiles with pdfLaTeX (Overleaf-ready). Standard packages only.
%=======================================================================
\documentclass[10pt,letterpaper]{article}

\usepackage[left=0.5in,right=0.5in,top=0.5in,bottom=0.5in]{geometry}
\usepackage[dvipsnames]{xcolor}
\usepackage{enumitem}
\usepackage{hyperref}
\usepackage{microtype}

% ---- colours ----------------------------------------------------------
\definecolor{accent}{HTML}{2A4B8D}   % deep analytical blue
\definecolor{muted}{HTML}{5A6473}

\hypersetup{colorlinks=true,urlcolor=accent,linkcolor=accent}

% ---- global layout ----------------------------------------------------
\pagestyle{empty}
\setlength{\parindent}{0pt}
\setlength{\tabcolsep}{0pt}
\renewcommand{\baselinestretch}{1.02}

% ---- helper macros ----------------------------------------------------
% Section header with coloured rule
\newcommand{\cvsec}[1]{%
  \vspace{7pt}%
  {\bfseries\large\color{accent}#1}\par
  \vspace{2pt}{\color{accent!55}\hrule height 0.7pt}\vspace{4pt}%
}

% Skill category: bold label on its own line, then the items
\newcommand{\skillcat}[2]{{\bfseries #1}\\ #2 \\[4pt]}

% Experience / project entry: title -- subtitle -- meta(date)
\newcommand{\entry}[3]{%
  \vspace{5pt}%
  {\bfseries #1}\hfill{\small\color{muted}#3}\\
  {\itshape #2}\par\vspace{1pt}%
}

% Tech-stack / tools line under an entry
\newcommand{\stack}[1]{{\footnotesize\color{muted}#1}\par\vspace{2pt}}

% Compact bullet list
\newenvironment{cvbullets}
  {\begin{itemize}[leftmargin=1.15em,itemsep=1.5pt,topsep=2pt,parsep=0pt,
     label=\textcolor{accent}{\footnotesize$\bullet$}]}
  {\end{itemize}}

%=======================================================================
\begin{document}

% ----------------------------- HEADER ----------------------------------
\begin{center}
  {\fontsize{22}{24}\selectfont\bfseries GABRIEL ADRIAN MANAS}\\[3pt]
  {\large\color{accent} Data Analyst \ \textbar\ \ Data Scientist}\\[5pt]
  {\small
    \href{mailto:a.gabrielmanas@gmail.com}{a.gabrielmanas@gmail.com}
    \ \textbar\ \ \href{mailto:gabrieladrian.manas@ust.edu.ph}{gabrieladrian.manas@ust.edu.ph}
    \ \textbar\ \ 0960 511 8496
    \ \textbar\ \ Manila / Cavite
    \ \textbar\ \ \href{https://linkedin.com/in/gabrielmanas}{LinkedIn}
  }
\end{center}
\vspace{2pt}

% ----------------------------- TWO COLUMNS -----------------------------
\noindent
\begin{minipage}[t]{0.31\textwidth}
%======================= LEFT COLUMN ===================================

\cvsec{Profile}
Medical physicist and quantitative researcher moving into \textbf{data
science}, with an MS in Applied Physics and a strong foundation in
\textbf{Python, statistics, and numerical modeling}. I turn measurement,
imaging, and simulation data into rigorous, decision-ready analysis --- and
communicate it through dashboards in \textbf{Power BI and Tableau}.

\cvsec{Technical Skills}
\skillcat{Programming}{Python, MATLAB, SQL}
\skillcat{Analysis \& Modeling}{pandas, NumPy, SciPy, statistical analysis, numerical methods, curve fitting, Monte Carlo (PHITS)}
\skillcat{Visualization \& BI}{Power BI, Tableau, Matplotlib, Excel}
\skillcat{Domain Data}{Imaging QA/QC, dosimetry, radiation measurement}
\skillcat{Tools}{Jupyter, Git, \LaTeX{}}

\cvsec{Education}
{\bfseries M.S. Applied Physics}\\
{\small (Medical Physics)}\\
University of Santo Tomas\\
{\small\color{muted}2026 -- Present}
\vspace{4pt}\\
{\bfseries B.S. Applied Physics}\\
{\small (Instrumentation)}\\
University of Santo Tomas\\
{\small\color{muted}2021 -- 2025}
\vspace{3pt}\\
{\footnotesize\itshape Thesis:} {\footnotesize Numerical Evaluation of
Bragg--Gray Cavity Integrals --- numerical modeling of stopping-power
ratios in radiation dosimetry (Python / MATLAB).}

\cvsec{Core Strengths}
\begin{cvbullets}
  \item Statistical \& numerical analysis
  \item Scientific computing in Python
  \item Data visualization \& storytelling
  \item QA/QC \& measurement data
\end{cvbullets}

\end{minipage}
\hfill
\begin{minipage}[t]{0.64\textwidth}
%======================= RIGHT COLUMN ==================================

\cvsec{Experience}

\entry{Junior Medical Physicist}{Medical Physics \& Health Physics Services, Inc. (MPHPS)}{Apr 2026 -- Present}
\stack{Python, Excel, Imaging QA}
\begin{cvbullets}
  \item Perform \textbf{performance testing and quality control} across
        imaging modalities (CT, radiography/fluoroscopy, mammography,
        PET/CT, SPECT/CT), capturing quantitative image-quality and
        radiation-output measurements.
  \item Analyze measurement data against regulatory tolerances to verify
        equipment performance and \textbf{flag deviations}, working with
        repeated QC metrics across machines and sites.
\end{cvbullets}

\entry{Science Teacher (Elementary -- Senior High)}{Hephzibah School Inc.}{Jun 2025 -- Apr 2026}
\begin{cvbullets}
  \item Taught science and quantitative reasoning, translating statistical
        and analytical concepts into clear, accessible explanations ---
        sharpening data-communication skills.
\end{cvbullets}

\entry{Medical Physics Intern}{UP -- Philippine General Hospital (PGH)}{May -- Jun 2024}
\stack{Excel, Dose Calculations, QA Documentation}
\begin{cvbullets}
  \item Collected, encoded, and validated \textbf{patient dose data} and
        performed activity and dose calculations across diagnostic
        radiology, nuclear medicine, and radiation oncology.
  \item Conducted radiation surveys and compiled QA/compliance
        documentation, organizing measurement datasets across X-ray, CT,
        fluoroscopy, MRI, SPECT/CT, PET/CT, brachytherapy, and LINAC systems.
\end{cvbullets}

\cvsec{Research \& Data Projects}

\entry{Undergraduate Thesis --- Bragg--Gray Cavity Integrals}{Python, MATLAB --- numerical analysis}{2025}
\begin{cvbullets}
  \item Built numerical models for stopping-power ratios in radiation
        dosimetry; \textbf{modeled mass electronic stopping power as a
        Gaussian} and evaluated cavity integrals by numerical integration.
\end{cvbullets}

\entry{Monte Carlo Simulation Study}{PHITS --- radiation transport}{2025}
\begin{cvbullets}
  \item Ran Monte Carlo simulations and analyzed simulation output to
        determine Bragg--Gray stopping-power ratios.
\end{cvbullets}

\entry{Research Dissemination}{Quantitative results, 3 venues}{2025}
\begin{cvbullets}
  \item Presented findings at the \textbf{43rd Samahang Pisika ng Pilipinas}
        Conference, the UST RCNAS Symposium, and the UST Mathematics \&
        Physics Seminar.
\end{cvbullets}

\end{minipage}

\end{document}
